\part{Knowledge, Attention, and Productive Competence}
\label{part:V}

Part V studies representations, memory, skill, automation, and verification as supports for, or threats to, competent continued action.

\begin{housing}Inspection records, contractor knowledge, and automated scheduling are examined, along with what a passed inspection establishes.\end{housing}

\chapter{Economic Representations and Semantic Compression}\label{ch:21}
Examines the representations used in planning, forecasting, classification, and coordination. A useful representation preserves distinctions relevant to action while compressing others. Defines task-relative representational adequacy.
\chspec{\Pref{P11}.}
\scaffold

\chapter{Memory as Organizational Capital}\label{ch:22}
Examines records, provenance, routines, and institutional memory as supports for continued competent action. Distinguishes retained information from retrievable and usable knowledge, and investigates when forgetting reduces costs without erasing obligations.
\chspec{\Pref{P21}.}
\draftnote{Candidate import: separate replay integrity from historical authority. A self-consistent checkpoint is not an independently trusted history, and branch selection is not evidence that the chosen branch is true. An institution can keep reproducible accounts without establishing that inputs are truthful or commitments legitimate.}
\scaffold

\chapter{Aspect Relegation and the Economics of Skill}\label{ch:23}
Models practice and routine as changes in the attention required to perform a task adequately. Relegation frees resources while creating dependence on stable conditions. Develops the economic significance of recognizing when a task must return to explicit supervision.
\chspec{\Pref{P16}.}
\draftnote{Candidate import: distinguish expenditure from damage to regenerative capacity, and allocate finite repair capacity through explicit response functions.}
\scaffold

\chapter{Automation and the Preservation of Intervention}\label{ch:24}
Assesses automation through output, error, oversight burden, retained competence, and recovery time. A formal example shows why removing human involvement can lower immediate costs while raising the cost of rare failures, with response time made endogenous to retained competence.
\chspec{\Pref{P16}.}
\draftnote{Candidate case (documentation-level only): an agent-automation tool. Isolated working copies are not an operating-system security boundary, and a command filter or stop control cannot recover effects that have already occurred. The economic question is how much review, intervention time, and recoverability the automation actually preserves. Exclude any zero-risk claim.}
\scaffold

\chapter{Verification, Proxies, and Institutional Confidence}\label{ch:25}
Distinguishes what an audit, benchmark, certification, or formal proof actually establishes. Examines the economic effects of treating a passed check as broader evidence than its specification warrants.
\chspec{\Pref{P17}.}
\draftnote{Candidate counterexample: a ledger verifier that checks seals against stored roots without recomputing them from contents, and checks debits one at a time without aggregating repeated postings, so that individually permissible debits jointly exceed the stated balance restriction. A check can succeed while omitting the relation a reader assumes it verifies.}
\scaffold
